\subsection{Elementary Row Operations}

\begin{definitionbox}[The three elementary row operations]
The elementary row operations are:
\begin{itemize}
    \item swapping two rows;
    \item multiplying one row by a nonzero scalar;
    \item adding a scalar multiple of one row to another row.
\end{itemize}
\end{definitionbox}

\begin{interpretationbox}[Why these operations are controlled]
Each operation can be reversed: a swap is undone by the same swap, a nonzero
scaling is undone by multiplying by the reciprocal, and adding a multiple of one
row is undone by subtracting the same multiple. This reversibility is why row
operations can transform a problem without losing its essential information.
\end{interpretationbox}

\begin{examplebox}[Eliminating information already carried by another row]
Consider
\[
A=
\begin{pmatrix}
1&2&-1\\
2&4&3\\
0&1&5
\end{pmatrix}.
\]
Replace the second row by the second row minus twice the first:
\[
(2,4,3)-2(1,2,-1)=(0,0,5).
\]
The transformed matrix is
\[
\begin{pmatrix}
1&2&-1\\
0&0&5\\
0&1&5
\end{pmatrix}.
\]
The first two entries of the old second row were already explained by twice the
first row. The new row isolates the part that was not.
\end{examplebox}

\begin{interpretationbox}[Three layers of a row operation]
A row operation contains arithmetic, a transformation of the matrix, and an
interpretation. The arithmetic produces the new row; the transformation changes
the representation; the interpretation explains what relation has become
visible.
\end{interpretationbox}
